\section{Related Work}

\subsection{Spacecraft Rendezvous and Proximity Operations}

Automated rendezvous and docking (ARD) has a multi-decade heritage
in human spaceflight and is now a core enabling capability for
emerging unmanned applications including on-orbit servicing, active
debris removal, and in-space assembly
\cite{fehse2003rendezvous,flores2014}.
The translational dynamics of close-proximity operations are
commonly described using the Hill--Clohessy--Wiltshire (HCW)
linearized equations, which provide a closed-form state-transition
solution for relative motion in the local orbital reference frame
under the assumption of a near-circular reference orbit
\cite{clohessy1960}.

Operational demonstration of autonomous rendezvous with a
non-cooperative target was achieved by the AVANTI experiment on
board the BIROS spacecraft, which validated relative navigation
and manoeuvring in low Earth orbit using optical measurements
alone \cite{gaias2018avanti}.
Such demonstrations highlight both the maturity of the underlying
guidance algorithms and the remaining challenge of obtaining
accurate pose estimates from passive camera measurements.

\subsection{Model-Based Relative Pose Estimation}

Recovering the six-degree-of-freedom (6-DoF) relative pose of a
spacecraft from monocular camera measurements is typically cast as a
Perspective-$n$-Point (PnP) problem, in which the camera pose is
recovered from known three-dimensional object points and their
two-dimensional image projections.
Kelsey~\textit{et al.}~\cite{kelsey2006} provided an early
demonstration that such vision-based pose estimation is viable
for autonomous rendezvous and docking applications.

Among PnP solvers, the EPnP algorithm of Lepetit~\textit{et al.}
\cite{lepetit2009epnp} is widely used in onboard applications
owing to its $\mathcal{O}(n)$ complexity: the reference points
are expressed as a weighted combination of four virtual control
points and the pose is recovered from a compact eigenvalue
problem.
A nonlinear Gauss--Newton step can subsequently minimize the
image reprojection error to further reduce the pose error, at
the cost of additional computation.

More recently, deep-learning methods have been applied to
non-cooperative spacecraft pose estimation, demonstrating improved
robustness to illumination variation and partial target visibility
\cite{sharma2018}.
However, these approaches require large labeled datasets and lack
the geometric guarantees of model-based pipelines.
The present work adopts the model-based approach, which is
interpretable, sample-efficient, and directly applicable when the
target geometry is known.

\subsection{Robust Correspondence Estimation}

Vision-based pose estimation degrades rapidly when keypoint
correspondences are incorrect, as can occur under partial
occlusion, background clutter, or mis-matched feature descriptors.
Random Sample Consensus (RANSAC), introduced by Fischler and
Bolles~\cite{fischler1981ransac}, provides a general framework
for robust estimation under outlier contamination.
RANSAC iteratively fits a minimal model hypothesis to a random
subset of correspondences and selects the hypothesis supported
by the largest consensus set of inliers.
The locally optimized RANSAC (LO-RANSAC) variant
\cite{chum2003loransac} additionally refits the model on the
full inlier set within each iteration, substantially improving
accuracy without increasing the required number of random
hypotheses.
The proposed framework uses RANSAC with EPnP as the hypothesis
solver and performs a final inlier refit to obtain the delivered
pose estimate.

\subsection{Recursive State Estimation}

Kalman-filter-based recursive estimators are the standard approach
for fusing noisy pose measurements with spacecraft dynamics models.
Because spacecraft attitude is parameterized by unit quaternions,
standard additive Extended Kalman Filters cannot preserve the
norm constraint through the update step.
The Multiplicative Extended Kalman Filter (MEKF) addresses this
by propagating a small rotation-vector error state that is
composed multiplicatively with the reference quaternion, ensuring
a unit-quaternion estimate at all times \cite{markley2003}.

The Unscented Kalman Filter (UKF), introduced by Julier and
Uhlmann~\cite{julier2004}, propagates a set of sigma points
through the nonlinear dynamics and avoids explicit Jacobian
computation.
The UKF is known to provide improved performance over the EKF
for strongly nonlinear systems, at the cost of a larger number
of function evaluations per step.
Both the MEKF and UKF are implemented and compared in the present
framework to characterize their relative performance under
identical controlled measurement conditions.

\subsection{Guidance and Control for Proximity Operations}

Linear Quadratic Regulation (LQR) provides an optimal closed-form
feedback law by solving the discrete algebraic Riccati equation
offline \cite{anderson1990}.
Under the HCW linearization, LQR yields a computationally
efficient infinite-horizon controller that is straightforward
to implement in the onboard flight computer.

Model Predictive Control (MPC) generalizes LQR by solving a
finite-horizon constrained optimization problem at each control
step, enabling explicit enforcement of thrust bounds and
trajectory constraints such as approach corridors
\cite{mayne2000}.
Breger and How~\cite{breger2008} demonstrated MPC-based safe
trajectory planning for spacecraft rendezvous, showing that
the receding-horizon framework can generate constraint-satisfying
manoeuvres in real time.
The present work evaluates both LQR and MPC for the same HCW
rendezvous scenario to quantify the trade-off between propellant
efficiency, terminal accuracy, and computational cost.

\subsection{Closed-Loop Simulation Frameworks}

Integrated numerical and physical testbeds have been developed
to evaluate complete navigation and control pipelines in a
controlled setting.
Gasbarri~\textit{et al.}~\cite{testbed2014} presented a
ground-based testbed for vision-based formation flying that
demonstrated camera-based relative navigation for multi-spacecraft
proximity operations.
The present work constructs a complementary numerical simulation
framework that integrates relative orbital dynamics, target
attitude propagation, synthetic camera measurements, pose
estimation, state estimation, and guidance and control within a
single closed-loop environment.
This modular architecture enables systematic evaluation of each
subsystem in isolation as well as full end-to-end performance
assessment under representative noise and outlier conditions.
